% Resumo – NBR 6028:2021 (parágrafo único, 150 a 500 palavras, voz ativa, 3.ª pessoa)
\begin{resumo}
Este trabalho investiga se empresas familiares brasileiras de capital aberto
apresentam endividamento diferente do das empresas não familiares. Com base na
perspectiva da riqueza socioemocional, a pesquisa parte da hipótese de que a
aversão à interferência de credores e à perda de controle leva as famílias
controladoras a adotarem políticas de financiamento mais conservadoras. A amostra
compreende 16 companhias do setor de saúde da B3, com dados trimestrais da
Economática de 2018 a 2025. As empresas foram classificadas como familiares
segundo o critério de \citeonline{santos2018}, e dois critérios alternativos foram
usados como testes de robustez. O endividamento foi medido por indicadores
contábeis de nível, de capacidade de pagamento e de composição da dívida, e as
diferenças entre os grupos foram avaliadas pelos testes t de Welch e de
Mann-Whitney, para o período completo e para os subperíodos anterior, durante e
posterior à pandemia de COVID-19. Os resultados indicam que as empresas
familiares apresentaram, em todos os exercícios, menor endividamento oneroso e
menor participação da dívida de curto prazo que as não familiares, mas essas
diferenças não foram estatisticamente significativas no período completo, e o
endividamento total dos dois grupos foi semelhante. A exceção é a composição da
dívida durante a pandemia, quando a menor participação da dívida de curto prazo
nas empresas familiares foi significativa ao nível de 5\%, em análise exploratória
por subperíodos, sem correção para múltiplos testes. A diferença no endividamento
oneroso inverteu o sinal sob os critérios alternativos de classificação e
mostrou-se sensível à exclusão de uma empresa sem dívida financeira. Conclui-se que
a hipótese não pode ser confirmada estatisticamente. A evidência descritiva é
compatível com uma postura mais conservadora das famílias em relação ao prazo da
dívida, mas também com diferenças de tamanho e de acesso ao crédito entre os
grupos.

\textbf{Palavras-chave}: empresas familiares; endividamento; estrutura de capital;
riqueza socioemocional; setor de saúde.
\end{resumo}
